\section{Engineering Discussion}
\label{sec:engineering-discussion}

The main engineering value of this project is the connection between analytical antenna theory,
full-wave CST simulation, and phased-array interpretation. The initial rectangular-patch dimensions
were estimated using the transmission-line model, whereas CST Studio Suite was required to account
for the finite substrate and ground-plane dimensions, practical feed geometry, fringing fields,
dielectric loss, port discontinuities, and mutual coupling between the array elements.

\subsection{Mesh-Limited CST Modeling Decision}

The CST model was initially prepared using finite-conductivity copper with
\[
\sigma_{\mathrm{Cu}}=5.87\times10^7~\mathrm{S/m}.
\]
However, the finite copper thickness of \(0.035~\mathrm{mm}\) introduced model dimensions that were
very small compared with the overall antenna size. A conventional volumetric conductor representation
produced a severe edge-length ratio and exceeded the CST Learning Edition mesh-cell limit during
refinement. The final baseline single-patch and array tuning runs therefore used perfect electric
conductor (PEC) metallization while retaining the DiClad 880-IM dielectric loss through
\[
\tan\delta=0.0009.
\]
This choice allowed the required resonance, matching, pattern, coupling, and beam-steering studies to
be completed within the available computational restriction.

After the baseline design was finalized, a supplementary finite-conductivity validation was performed
for the single patch using CST's \emph{Lossy metal} model. An adaptive-mesh attempt again exceeded the
\(20{,}000\)-cell limit before convergence and was discarded. A fixed-mesh run with the same geometry
completed successfully and produced a resonance near \(2.504~\mathrm{GHz}\),
\(S_{11,\min}\approx-39.72~\mathrm{dB}\), and directivity near \(7.23~\mathrm{dBi}\). These values
were nearly identical to the PEC baseline. The radiation and total efficiencies decreased to
approximately \(92.46\%\) and \(92.38\%\), respectively, compared with \(96.3\%\) and \(96.2\%\)
for PEC. The detailed comparison is provided in Section~\ref{subsec:single-patch-copper-validation}.

The validation demonstrates that PEC was an effective approximation for resonance, input matching,
directivity, and normalized pattern shape in the present design. Its principal limitation was an
overestimate of efficiency because finite copper ohmic loss was omitted. The fixed-mesh copper result
should not be interpreted as a formal mesh-converged solution, but it provides a useful conductor-model
sensitivity check within the Learning Edition limit.

\subsection{Single-Patch Tuning Behavior}

The single-patch tuning process demonstrated a physically consistent relationship between patch length
and resonant frequency. The initial analytical patch length of \(39.657~\mathrm{mm}\) produced a
resonance near \(2.431~\mathrm{GHz}\), below the selected design frequency of
\(2.501~\mathrm{GHz}\). Because the patch was electrically too long, the patch length was reduced in
successive tuning runs.

The final single-patch design used
\[
L=38.40~\mathrm{mm},
\qquad
y_{\mathrm{inset}}=14.11~\mathrm{mm},
\]
and produced
\[
S_{11,\min}\approx -39.87~\mathrm{dB}
\]
at approximately
\[
f_{\mathrm{r}}=2.50507~\mathrm{GHz}.
\]
The resonant-frequency error relative to the selected design target was therefore approximately \(0.16\%\).

The simulated single-patch farfield exhibited the expected broadside radiation behavior of a
rectangular microstrip patch. The PEC baseline peak directivity was approximately
\(7.226~\mathrm{dBi}\), with radiation and total efficiencies of approximately \(96.3\%\)
and \(96.2\%\), respectively. The supplementary finite-conductivity run preserved essentially the
same resonance, matching, directivity, and normalized pattern while reducing the radiation and total
efficiencies to approximately \(92.46\%\) and \(92.38\%\). This behavior is physically consistent
with the addition of copper ohmic loss.

\subsection{Single Patch Versus Array Tuning}

When the tuned single patch was duplicated to form the \(2\times1\) array, the resonance shifted
substantially downward. The initial array configuration, which retained the final isolated-patch
dimensions, produced resonance near \(2.399~\mathrm{GHz}\).

This shift occurred because the individual elements no longer radiated in isolation. Mutual coupling
altered the current distribution, input impedance, and effective electrical length of each patch.
In addition, the enlarged common substrate and ground-plane environment modified the electromagnetic
loading of the elements. Consequently, the optimized dimensions of the isolated patch could not be
used directly as the final dimensions of the array.

The array was therefore retuned independently. The final array geometry used
\[
L=36.70~\mathrm{mm},
\qquad
y_{\mathrm{inset}}=13.35~\mathrm{mm},
\qquad
d=59.935~\mathrm{mm}.
\]
The selected spacing corresponds approximately to
\[
d\approx\frac{\lambda_0}{2}
\]
at the selected operating frequency.

Near the target frequency, the final tuned array produced approximately
\[
S_{11}\approx -11.9~\mathrm{dB},
\qquad
S_{22}\approx -15.2~\mathrm{dB}.
\]
The mutual coupling was approximately
\[
S_{12}\approx S_{21}\approx -15~\mathrm{dB}.
\]

The final array was selected as an engineering compromise among resonant-frequency proximity,
acceptable matching at both input ports, and manageable mutual coupling. The coupling result confirms
that the two patches are not electromagnetically independent. Excitation of one element produces
fields and induced currents on the neighboring element, thereby affecting the input match, radiation
behavior, and phased-array response.

\subsection{Beam-Steering Interpretation}

Beam steering was implemented using CST amplitude-phase post-processing. The independently simulated
farfields associated with the two ports were combined using equal amplitudes and different relative
phases.

The three evaluated excitation cases were
\[
\left(1\angle0^\circ,\;1\angle0^\circ\right),
\]
\[
\left(1\angle0^\circ,\;1\angle-90^\circ\right),
\]
and
\[
\left(1\angle0^\circ,\;1\angle+90^\circ\right).
\]

The \(0^\circ\) relative-phase case produced the broadside reference pattern, with the two principal
cuts locating the main beam at approximately \(0^\circ\) and \(1^\circ\). Applying a relative phase
shift of \(-90^\circ\) displaced the main lobe by approximately \(20^\circ\) toward the
\(\phi=0^\circ\) half-plane, whereas reversing the relative phase to \(+90^\circ\) produced an
approximately \(20^\circ\) displacement toward the \(\phi=180^\circ\) half-plane.

For a two-element array with approximately half-wavelength spacing, the ideal array-factor relationship is
\[
\Delta\phi=-k_0d\sin\theta_0.
\]
With
\[
d\approx\frac{\lambda_0}{2},
\]
this becomes
\[
\Delta\phi\approx-180^\circ\sin\theta_0.
\]
Therefore, relative phase shifts of \(-90^\circ\) and \(+90^\circ\) ideally correspond to equal scan-angle
magnitudes of approximately \(30^\circ\) in opposite directions. The CST-observed magnitude of
approximately \(20^\circ\) is smaller than the ideal prediction.

Exact agreement between the ideal array-factor prediction and the CST farfield is not necessarily
expected. The analytical array factor assumes idealized point sources, whereas the CST model includes
the finite patch element pattern, substrate loading, feed geometry, port definition, finite ground
plane, and mutual coupling. In addition, the displayed sign of the steering direction depends on the
CST coordinate system and the ordering of Ports 1 and 2.

The CST post-processing results nevertheless demonstrated the essential phased-array mechanism:
the direction of the combined main beam was controlled by the relative excitation phase while the
excitation amplitudes remained equal.

CST reported that unnormalized balance information was unavailable for the combined post-processed
farfield results. Consequently, the combined-result gain and efficiency values were not treated as
validated quantitative metrics. CST displayed peak directivities of approximately
\(9.655~\mathrm{dBi}\), \(9.489~\mathrm{dBi}\), and \(9.531~\mathrm{dBi}\) for the broadside,
\(-90^\circ\), and \(+90^\circ\) cases, respectively. The combined farfields and principal-plane cuts
were used primarily to evaluate directivity-pattern shape, beamwidth, and the direction of beam
steering.

\subsection{Numerical Reliability and Convergence Considerations}
\label{subsec:numerical-reliability}

The interpretation of the CST results must account for the numerical limitations of the available
simulation environment. The PEC baseline tuning runs and the successful finite-conductivity validation
used a fixed tetrahedral mesh so that the models remained within the CST Learning Edition cell limit.
A formal mesh-convergence study was therefore not completed.

A separate adaptive-mesh finite-conductivity run was attempted as a numerical reliability check. The
solver completed several adaptive passes, but the mesh exceeded the \(20{,}000\)-cell license limit
before convergence was reached. No quantitative result from that incomplete run was used in the
report. The successful lossy-metal result was obtained only after mesh adaptation was disabled and is
therefore identified explicitly as a fixed-mesh conductor-model sensitivity study.

A formal boundary-convergence study was also not performed. The final simulations used open
boundaries with added space, but the boundary distance was not systematically increased and compared
across multiple runs. Accordingly, the reported results should be interpreted as engineering-level
simulation results rather than fully converged production-level predictions.

Despite these limitations, the principal physical trends were internally consistent. For the isolated
single patch, reducing the patch length shifted the resonance upward toward the target frequency. For
the array, adding the second element caused a downward resonance shift accompanied by quantifiable
mutual coupling. The separate PEC and finite-conductivity single-patch models produced nearly
identical resonance, matching, directivity, and normalized pattern shape, while the finite-conductivity
model produced the expected reduction in efficiency. No unexplained non-monotonic trend was used as
the basis for the final design.

A more rigorous numerical validation procedure would repeat the final single-patch and array models
using several progressively refined local meshes. For each mesh level, the following quantities should
be compared:
\begin{itemize}
    \item resonant frequency;
    \item minimum \(S_{11}\) and \(S_{22}\);
    \item mutual coupling \(S_{12}\) and \(S_{21}\);
    \item peak directivity;
    \item radiation and total efficiencies; and
    \item beam-steering angle.
\end{itemize}

Boundary convergence could be evaluated by repeating the simulations with progressively increased
open-boundary spacing and confirming that the changes in resonant frequency, \(S\)-parameters, and
farfield quantities become negligible. A higher-capacity solver license would also permit adaptive
finite-conductivity validation of the complete \(2\times1\) array.

Within the computational limitations of the CST Learning Edition, the analytical calculations,
single-element tuning, conductor-model sensitivity check, two-port array characterization,
mutual-coupling results, and phased farfield combinations were mutually consistent. The final model
therefore provides a credible demonstration of the complete engineering workflow from analytical patch
design to phase-controlled array beam steering.
